% This is samplepaper.tex, a sample chapter demonstrating the
% LLNCS macro package for Springer Computer Science proceedings;
% Version 2.21 of 2022/01/12
%
\documentclass[runningheads]{llncs}
%
\usepackage[T1]{fontenc}
% T1 fonts will be used to generate the final print and online PDFs,
% so please use T1 fonts in your manuscript whenever possible.
% Other font encondings may result in incorrect characters.
%
\usepackage{graphicx}
\usepackage{hyperref}
\usepackage{listings}
\usepackage{xcolor}

\lstdefinelanguage{Solidity}{
	morekeywords={pragma, solidity, contract, struct, mapping, 
		address, uint256, bytes32, string, bool, 
		function, public, private, view, returns, 
		emit, event, require, memory, storage},
	sensitive=true,
	morecomment=[l]{//},
	morecomment=[s]{/*}{*/},
	morestring=[b]{"}
}

\lstset{
	basicstyle=\ttfamily\small,
	keywordstyle=\color{blue}\bfseries,
	commentstyle=\color{gray}\itshape,
	stringstyle=\color{orange},
	breaklines=true,
	showstringspaces=false,
	frame=single
}

\renewcommand{\theHsection}{\thepart.\thesection} 



\begin{document}
	%
	\title{EcoPassEU: A Privacy-Aware Digital Product Passport DApp for Sustainable Consumer Goods in Europe}
	
	\titlerunning{EcoPassEU: A Privacy-Aware Digital Product Passport DApp}
	
	\author{Fuyu Cao\inst{1} \and
		Luxiao Cao\inst{1} \and
		Detai Dong\inst{1} \and
		Akshansh Rajora\inst{1} \and
		Tianwei Yu\inst{1}}
	
	\authorrunning{F. Cao et al.}
	
	\institute{University of Bristol, Bristol, UK \\
		\email{\{1d25515, ne25392, mz25325, tk25240, np25024\}@bristol.ac.uk}}
	%
	\maketitle              % typeset the header of the contribution
	%

	%
	%
	\section{Introduction}
	Every year, billions of products enter the European market with  sustainability labels, yet consumers have little reliable way of checking whether claims such as ``low carbon'' or ``ethically sourced'' are accurate. A study commissioned by the European Commission found that more than half of the green environmental declarations lack verifiable evidence to support, eroding public confidence in corporate environmental commitments \cite{Silva2021The}.
	
	To address this credibility gap, the Ecodesign for Sustainable Products Regulation (ESPR) now mandates Digital Product Passports (DPPs) —standardised lifecycle records—for most EU products \cite{circularise_dpp_2025}.  
	With the implementation of compliance in 2027, reliable DPP solutions have become a market urgent need\cite{circularise_dpp_2025}.
	
	Our product EcoPassEU is a blockchain-based DPP platform which  allows manufacturers to record product information as permanent, publicly verifiable entries on the Ethereum blockchain. 
	Unlike traditional databases, the immutability of blockchain records means that once the data is released, it cannot be altered retrospectively, eliminating the possibility of post-hoc falsification \cite{Casino2019Systematic}. 
	Any user can obtain a complete and encrypted product history by entering the product ID or scanning the QR code.
	
	EcoPassEU serves four key stakeholders. 
	It provides SME manufacturers with a cost-effective ESPR compliance channel through a web for batch data entry and automated QR code generation. 
	At the same time, consumers, regulators and retailers can also get open and wallet-free access. 
	Just search for the product ID to immediately obtain the blockchain-verified product passport and supply chain map, enabling the independent verification of environmental claims without prior technical expertise.
	
	The value of EcoPassEU is threefold. 
	Primarily, it provides manufacturers with a practical way to ESPR deadlines, reducing the risk of market exclusion. 
	Meanwhile, the platform directly cracked down on greenwashing, 
	securing product records via cryptographic hashing replaces brand self-reporting with independent, third-party verification.
	Thirdly, it promotes the development of circular economy by equipping recyclers with the precise material data necessary to recover resources when products are scrapped.
	
	
	
	
	
	
	\section{Background}
	\subsection{Digital Product Passports and the Regulatory Context}
	A Digital Product Passport (DPP) is a structured digital record encompassing lifecycle data—such as material composition, provenance, and end-of-life guidance—designed to inform stakeholders and support more sustainable design choices, efficient resource use, and the transition towards circular business models
	\cite{Adisorn2021Towards,Basal2024Digital,VanCapelleveen2023The,Jansen2023Stop}.
	
	
	The regulatory foundation for DPPs comes from the European Green Deal and Circular Economy Action Plan \cite{eu_green_deal_2019,ceap_2020}. Based on these 2 frameworks, the Ecodesign for Sustainable Products Regulation (ESPR), came into force on 18 July 2024,, mandates DPPs for almost all physical goods sold in the EU, regardless of the seller's location \cite{circularise_dpp_2025}.
	
	\subsection{The Problem: Trust in Sustainability Claims}
	Despite regulatory momentum,  the implementations of most DPPs rely on centralised, proprietary databases \cite{Psarommatis2024Digital}. 
	This architecture encourages the selective concealment of data, hindering aftermarket participants like independent repairers \cite{Ducuing2023Data}. 
	Furthermore, managing huge volume of lifecycle data centrally raises significant concerns regarding storage efficiency, data integrity, and vulnerability to breaches \cite{Voulgaridis2024Digital}. 
	Taken together, downstream actors — including importers, recyclers, and end consumers — lack reliable ways to independently verify product information.
	
	This problem is exacerbated by the prevalence of greenwashing, broadly understood as false, vague, or unsubstantiated environmental claims that make products appear more sustainable than they are \cite{Kundi2024The}. 
	A 2020--2021 EU sweep of 344 online claims  found that nearly half were potentially deceptive—often substituting verifiable details with vague terms like ``conscious''—and 59\% lacked accessible evidence \cite{Silva2021The}. 
	Another EU website screening found that 42\% of sustainability claims were misleading enough to constitute improper commercial behaviour \cite{Silva2021The}. 

	
	Blockchain provides a structural solution to this trust problem. 
	Because records on a public blockchain are replicated across a distributed network, neither single party can unilaterally tamper with the records, providing a form of tamper evidence that cannot be replicated by a centralised system \cite{Casino2019Systematic,Nakamoto2008Bitcoin}.
	Smart contract introduced in the Ethereum protocol allows data access and update rules to be encoded in publicly readable logic, so that there is no need for trusted intermediaries and enabling cross-organisational verification without shared infrastructure \cite{Wood2014Ethereum}.
	
	
	\subsection{Related Work and Existing Platforms}
	Blockchain-based product traceability uses its immutability and decentralised verification to improve supply chain transparency \cite{Casino2019Systematic}. 
	Blockchain-based product traceability emerged as a research topic around 2015--2016, with early work concentrated in food and agriculture before spreading to pharmaceuticals, electronics, and textiles \cite{Ellahi2023Blockchain-Based,Tian2016Agri}. 
	More recently, research has shifted from feasibility towards scale and efficiency: for instance, Walmart's use of IBM Food Trust — built on Hyperledger Fabric — to reduce food-item tracing time from six days to under three seconds \cite{Ellahi2023Blockchain-Based} .
	
	In the commercial field, Circularise is among the closest comparable 
	platforms, offering DPPs with selective disclosure of sustainability data 
	across supply chains, with applications spanning the automotive, battery, and textile sectors \cite{Circularise2026Platform}.
	VeChain similarly deploys blockchain and smart contracts on its VeChainThor platform to achieve supply chain transparency, though it runs on on a permissioned network with a restricted set of authority-approved validators \cite{VeChain2026}
	Both platforms have a common limitation: reliance on proprietary infrastructure or permissioned networks introduces vendor dependency and reduces interoperability \cite{Belchior2020A,Wang2023Exploring}.
	
	EcoPassEU takes a different approach. 
	Rather than relying on a proprietary database, it is built on the public Ethereum blockchain and IPFS, with all rules encoded in open smart contracts. 
	Product data is stored on IPFS, whilst only a small set of cryptographic identifiers is written on-chain, keeping costs low. 
	The batch registration function further reduces the cost per product, making the platform viable for SMEs. 
	And as an open-source project, it carries no vendor lock-in risk.
	
	
	\subsection{Business Viability}
	EcoPassEU's revenue model rests on two foundations.
	The first is public funding: the EU's CIRPASS-2 project (2024--2027) explicitly subsidise DPP deployment across textiles, electronics, and construction, with dedicated support for SMEs adopting 
	``DPP-as-a-Service'' models \cite{cirpass2_2024}. 
	As an open source platform, EcoPassEU is  fully capable of benefitting from such grants.
	
	Another is subscription revenue. 
	ESPR compliance is legally mandatory for any business selling physical goods in the EU, forming a captive market independent of voluntary commitments \cite{circularise_dpp_2025}
	The green technology and sustainability market is projected to grow from USD~25.47 billion in 2025 to USD~73.90 billion by 2030 \cite{marketsandmarkets2025}.
	EcoPassEU charges a graded subscription fee based on passport volume, with optional premium tiers for API integration and compliance reporting. 
	The platform is built on public Ethereum, supporting payment with ETH or other ERC-20 tokens, providing international users a lower-friction settlement option and removing the need for traditional banking intermediaries.
	
	
	
	
	
	\section{Design Documentation}
	\subsection{Design Philosophy and Architectural Choices}
	EcoPassEU is a DPP designed to deliver ESPR-compliant Digital Product Passports with independently verifiable sustainability data. The design is governed by three design principles, each addressing a specific product requirement. 
	
	First, a hybrid storage model is adopted: full product data is pinned to 
	IPFS via Pinata \cite{Benet2014IPFS}, whilst only the content identifier 
	(CID), a \texttt{keccak256} integrity hash, the issuer address, and a 
	timestamp are recorded on-chain \cite{Wood2014Ethereum}, keeping gas costs 
	low. Second, commercially sensitive data remains off-chain by design; the 
	on-chain hash serves as an integrity anchor, confirming that the off-chain 
	document is unaltered without exposing its contents to the ledger. Third, 
	all read queries use a public RPC endpoint, allowing consumers and 
	regulators to inspect any passport without holding a wallet or incurring 
	gas fees.

	\subsection{System Architecture}
	Figure ~\ref{fig:architecture} shows the four-layer architecture of the overall system
    \vspace{-0.2cm}
	\begin{figure}[h]
		\centering
		\includegraphics[width=\linewidth]{figures/1.png}
		\caption{EcoPassEU system architecture.}
		\label{fig:architecture}
	\end{figure}
    \vspace{-0.2cm}
	
	At the user layer, two distinct actors participants interact with the system: manufacturers, who connect through MetaMask to sign and submit transactions, and consumers or regulators, who can query product records without a wallet. 
	Both routes pass through the application layer — a React/Vitefrontend deployed on Vercel, using Ethers.js~v6 for contract calls — which is responsible for uploading documents to IPFS, calling the Sepolia test network and local integrity check. 
	The \texttt{PassportRegistry} smart contract records the CID, metadata hash, issuer address, and timestamp on the Sepolia ledger, emitting a \texttt{PassportIssued} event upon registration. 
	An \texttt{exists} flag can prevent duplicate entries. Supply chain provenance is visualised via Leaflet.js with OpenStreetMap.
	
	\subsection{Use Cases and Workflow}
	EcoPassEU serves three actor types with differing technical needs, as illustrated in Figure~\ref{fig:usecase}. 
	\vspace{-0.2cm}
	\begin{figure}[h]
		\centering
		\includegraphics[width=0.75\linewidth]{figures/3.png}
		\caption{EcoPassEU use case diagram.}
		\label{fig:usecase}
	\end{figure}
	\vspace{-0.2cm}
	
	Wallet connection via MetaMask on the Sepolia testnet (UC-01) is a common premise for all manufacturers and administrators, as each on-chain write must be signed by a wallet key pair and cryptographically linked to an accountable issuer \cite{Jorgensen2022Universal,Keršič2023Orchestrating}. To register a product, the manufacturer fills in a structured registration form (UC-02), including product details, material composition, and supply chain sources. Upon submission, the frontend serialises the form data to JSON and uploads it to IPFS via Pinata (UC-03) to receive a CID. 
	
	Next, the system calculates the \texttt{keccak256} hash of the JSON locally and passes it to the \texttt{registerPassport()} function on the smart contract alongside the CID and product ID (UC-04). After the user confirms the transaction in MetaMask, a \texttt{PassportIssued} event is emitted, and a QR code with an Etherscan link is provided for users to download (UC-05), which also accesses the product list view (UC-06). The admin role requires the same wallet connection but adds batch registration capabilities (UC-12) to reduce the unit gas cost for high-volume producers.
	
	On the read side, consumers and regulators require no wallet to verify a product. The user submits a product ID (UC-07), and the frontend queries the contract via a public RPC endpoint to jump to the full passport view (UC-08). If the passport exists, the system obtains the complete JSON data from IPFS and recalculates its hash on the client: a green integrity badge is displayed if it matches, and a warning if it does not match (UC-09). The product passport page will also display an interactive supply chain map (UC-10) and an optional issuer profile (UC-11). This separation between the write-gated and read-open flows supports the design goal of providing transparent information without setting financial or technical barriers to end users.
	
	
	\subsection{Sequence and State Diagrams}
	
	Figure ~\ref{fig:sequence} outlines the technical details behind the core application flows, focussing on the strict execution order required for successful registration.
	Specifically, the frontend must first obtain a valid CID from IPFS before prompting the wallet signature and triggering the smart contract. 
	This ensures that the on-chain pointer never leads to data loss. 
	Meanwhile, the verification flow demonstrates how the system relies entirely on the client logic to compare the acquired IPFS documents against the on-chain hash, ensuring trustless validation.
	
	\vspace{-0.2cm}
	\begin{figure}[htbp]
		\centering
		\includegraphics[width=0.75\linewidth]{figures/2.png}
		\caption{UML sequence diagram detailing the registration and wallet-free verification flows.}
		\label{fig:sequence}
	\end{figure}
	\vspace{-0.2cm}
	
	Figure~\ref{fig:state} presents an UML state machine diagram to track the digital passport's lifecycle. 
	Physical products initially entered the system as \texttt{Unregistered}.
	Only after the data is successfully uploaded to Pinata will the status change to \texttt{IPFSStored}.
	Following the user's signature, the status changes to \texttt{TxPending}, and then after mining the block on Sepolia, it reaches the immutable \texttt{Registered}.
	During the final verification stage, the client hash comparison will mark the record as either \texttt{Verified} (if the hash matches) or \texttt{Compromised} (if the data has been altered).
	
	\vspace{-0.2cm}
	\begin{figure}[htbp]
		\centering
		\includegraphics[width=0.4\linewidth]{figures/4.png}
		\caption{UML state machine diagram showing the lifecycle of a 
			digital product passport.}
		\label{fig:state}
	\end{figure}
	\vspace{-0.2cm}
	
	\subsection{User Interface Design}
	
	\vspace{-0.2cm}
	\begin{figure}[htbp]
		\centering
		\includegraphics[width=0.9\linewidth]{figures/5.png}
		\caption{Wireframe of the search interface, showing the initial state and the on-chain result card.}
		\label{fig:ui_search}
	\end{figure}
	\vspace{-0.2cm}
	
	
	
	To reflect the environmental theme, the page adopts a dark green colour scheme (\texttt{\#2d6a4f}, \texttt{\#1b4332}). The layout uses standard \texttt{rem} spacing for a clean and consistent visual effect. As shown in Figure~\ref{fig:ui_search}, the search interface is straightforward. Users can query a product ID via the central search bar or use the predefined quick search buttons, entirely without connecting a wallet. The interface transitions from an initial empty state to a verified result card upon a successful search. This card immediately displays the core on-chain data—such as the issuer address, registration timestamp, and metadata hash—providing instant verification before the user proceeds to view the full product passport.
	

	
	
	Figure~\ref{fig:ui_register} displays the registration layout. 
	A simple tracker (Fill Details $\to$ Upload IPFS $\to$ Register Chain $\to$ Done) shows the user's registration progress. 
	All form fields—product details, materials, and supply chain—are displayed on a scrollable page instead of hiding sections.
	If MetaMask is not connected, an amber warning banner will be displayed to stop users from uploading data until they log in.
	
	\vspace{-0.2cm}
	\begin{figure}[htbp]
		\centering
		\includegraphics[width=\linewidth]{figures/6.png}
		\caption{Wireframe layout of the manufacturer registration form.}
		\label{fig:ui_register}
	\end{figure}
	\vspace{-0.2cm}
	
	As seen in Figure~\ref{fig:ui_verify}, the result page clearly divides the information. 
	The \emph{Blockchain Record} card displays on-chain data (issuer, time, and hash), while other cards display IPFS details. 
	Integrity badge—a green check for a match, or an amber triangle for an error—makes it easy to judge whether the data is valid. 
	If a user does not install MetaMask, the page will automatically use public RPC to load information.
	
	\vspace{-0.2cm}
	\begin{figure}[htbp]
		\centering
		\includegraphics[width=\linewidth]{figures/7.png}
		\caption{Wireframe detailing the structure of the product verification page.}
		\label{fig:ui_verify}
	\end{figure}
	\vspace{-0.2cm}
	
	Our site is built with CSS Grid and Flexbox to run perfectly on any device. 
	At 600\,px, forms will switch to a single-column layout to adapt to the mobile device screen. 
	A 768\,px breakpoint adjusts the menu and interactive map sizes for tablet users.
	
	
	
	
	
	
	\section{Project Execution}
	
	\subsection{Core Code Analysis and Data Structures}
	EcoPassEU uses a hybrid architecture, separating on-chain and off-chain data to balance transparency with efficiency. The core data structure is divided across these two environments. 
	
	Off-chain, the complete lifecycle data---including category, brand, material composition, and supply chain provenance---is structured as JSON objects and pinned to IPFS via Pinata. On-chain, the \texttt{PassportRegistry.sol} smart contract uses a strictly typed \texttt{Passport} struct to store only the essential, immutable references:
	
	\begin{lstlisting}[language=Solidity, 
		caption={Core Data Structure in PassportRegistry.sol}]
		struct Passport {
			string  ipfsCID;       // off-chain document pointer
			bytes32 metadataHash;  // keccak256 of JSON for tamper-evidence
			address issuer;
			uint256 timestamp;
			bool    exists;
		}
		
		// productId => Passport
		mapping(string => Passport) private passports;
	\end{lstlisting}
	
	This private mapping directly associates a unique product ID with its corresponding 
	passport data for O(1) retrieval. Initially, we considered using ERC-1155 tokens 
	to represent these passports; however, we opted for this lightweight registry 
	approach. Storing only the IPFS CID, a cryptographic hash, the issuer address, 
	and a block timestamp on-chain significantly reduces deployment complexity and 
	gas costs per registration, whilst keeping the full product lifecycle data 
	off-chain on IPFS.
	
	\subsection{Testing Strategy}
	To ensure the reliability of the smart contracts, we implemented a comprehensive unit testing suite using the Hardhat environment alongside the Chai assertion library. This follows established practices for smart contract validation \cite{Ferreira2020SmartBugs}. 
	
	Our testing strategy focused on core functionality and edge cases. We confirmed that the \texttt{registerPassport} function accurately stores the transaction sender's address as the issuer, preventing forgery, and generates a valid block timestamp. Furthermore, boundary testing verified that the contract successfully reverts transactions and reports errors when a duplicate product ID is submitted. We also tested the \texttt{batchRegisterPassports} function to ensure it can process arrays efficiently, emit correct events for valid entries, and silently skip already registered products without failing the entire batch.
	
	
	\subsection{Final UI Showcase}
	
	
    Figures~\ref{fig:ui_register} and~\ref{fig:ui_verify} show the 
    deployed EcoPassEU interface.
	The registration portal guides manufacturers through a structured 
	form with a step indicator, whilst the public verification page 
	displays the passport details, integrity status, and IPFS document 
	link without requiring a wallet.
	\vspace{-0.2cm}
	\begin{figure}[htbp]
		\centering
		\includegraphics[width=0.85\linewidth]{figures/8.png}
		\caption{Deployed EcoPassEU interface: manufacturer registration portal.}
		\label{fig:ui_register}
	\end{figure}
	
	\begin{figure}[htbp]
		\centering
		\includegraphics[width=0.85\linewidth]{figures/9.png}
		\caption{Deployed EcoPassEU interface: public product verification page.}
		\label{fig:ui_verify}
	\end{figure}
	\vspace{-0.2cm}
	
	\subsection{User Manual}
	The application is designed to be intuitive for two distinct user groups. 
	
	\textbf{For Manufacturers (Registration):}
	\begin{enumerate}
		\item Connect your MetaMask wallet to the Sepolia testnet using the top navigation bar.
		\item Navigate to the 'Register' page and complete the product detail forms, including material and supply chain data.
		\item Click 'Upload to IPFS'. The system will generate a unique CID.
		\item Click 'Register on Blockchain' and approve the transaction in the MetaMask pop-up.
		\item Once confirmed, download the generated QR code to attach to your physical products.
	\end{enumerate}
	
	\textbf{For Consumers and Regulators (Verification):}
	\begin{enumerate}
		\item No wallet connection is required. Open the EcoPassEU homepage.
		\item Enter the Product ID into the central search bar, or scan the product's QR code.
		\item Review the product's lifecycle data. A green 'Blockchain Verified' badge indicates that the IPFS data matches the immutable blockchain record.
	\end{enumerate}
	
	
	
	
	
	
	
	
	
	
	
	
	
	
	
	\subsection{Critical Evaluation}
	
	A critical assessment of the execution cost of the DApp shows that its gas efficiency has been highly optimised. The design of storing larger string data such as repair guides or recycling instructions off-chain ensures that the state changes required for the registry are extremely lightweight. According to the testnet deployment results and standard EVM opcodes estimates, executing a single registration transaction consumes about 125,430 gas. In addition, the \texttt{batchRegisterPassports} function shows significant operational scalability. By traversing a set of ID and hash values in one transaction, the contract avoids the basic transaction fee of 21,000 gas. Test data shows that when 10 products are registered in bulk, the average gas cost of a single product can be reduced to about 92,100 gas, achieving an efficiency improvement of more than 25\%, making the platform economically feasible for small and medium-sized enterprises.
	
	In terms of security, the implementation has successfully reduced the risk of tampering and unauthorised modification, which is a central concern in Ethereum smart contract security \cite{Luu2016Making,Tsankov2018Securify}. In smart contracts, the explicit \texttt{require} check provides a strict tamperability detection guarantee by making the registered passport immutable. In addition, the stored metadata hash establishes a cryptographic commitment between the on-chain registry and the off-chain IPFS load. If the malicious actor tries to intercept and tamper with the IPFS document, the cryptographic hash of the tampered document will no longer match the immutable record saved on the chain, thus immediately invalidating the passport. Unlike existing solutions that rely on permissioned nodes or proprietary databases, EcoPassEU achieves enterprise-level privacy protection on a fully open network through the IPFS hash pairing mechanism.
	
	Although the decentralised architecture provides strong security, it also brings some runtime restrictions. The DApp performs quickly in on-chain verification and can usually be read in the state of completion in milliseconds. However, the access to the complete product history depends on public IPFS gateways, which sometimes have higher latency than centralised cloud servers. In addition, although the current platform runs efficiently on the Sepolia testnet, if it is migrated to the Ethereum mainnet, users will face volatile gas costs. Future versions of EcoPassEU should be considered for deployment on Layer-2 networks such as Polygon. This can not only maintain EVM compatibility, but also achieve near-instant final confirmation; at the same time, if it is necessary to support secondary market transfer in the future, it can be further reconstructed into a plan based on ERC-1155 tokens \cite{Radomski2018ERC1155}.
	
	
	\section{Conclusion}
	% Concise overview of findings, achievements, limitations, and future work.
	Summarize the main theoretical and technical findings. Detail what the group achieved, acknowledge limitations (e.g., the physical-to-digital bridge vulnerabilities), and propose future work.
	
	
	
	\newpage
	%
	% ---- Bibliography ----
	%
	% BibTeX users should specify bibliography style 'splncs04'.
	% References will then be sorted and formatted in the correct style.
	%
	\bibliographystyle{splncs04}
	\bibliography{mybibliography}
	
	\newpage
	\appendix
	\section{Appendix}
	% The appendix has no page limit and must include specific items in order.
	This appendix includes the mandatory materials required for the FTGP summative assessment.
	
	\subsection{Smart Contract Deployment}
	\begin{itemize}
		\item \textbf{Smart Contract Address(es):} [Insert Sepolia testnet addresses here]
		\item \textbf{Etherscan Links:} [Insert links to verified contracts on Etherscan]
	\end{itemize}
	
	\subsection{Project Links}
	\begin{itemize}
		\item \textbf{Deployed Website:} [Insert URL]
		\item \textbf{GitHub Repository:} \url{https://github.com/DeTaiDong/FTGP2425_Group1_11} 
	\end{itemize}
	
	\subsection{Transaction Hashes}
	% Any deploy/test transaction hashes along with brief explanations.
	\begin{itemize}
		\item \textbf{Deployment TX:} [Insert Hash] - Explanation.
	\end{itemize}
	
	\subsection{Equity Shares \& Contributions}
	% Distribution of equity shares and screenshot of GitHub Insights
	\begin{itemize}
		\item Luxiao Cao: [Weight]
		\item Akshansh Rajora: [Weight]
		\item Fuyu Cao: [Weight]
		\item Tianwei Yu: [Weight]
		\item Detai Dong: [Weight]
	\end{itemize}
	\noindent \textit{[Insert Screenshot of GitHub - Insights - Contributor here]}
	
	\subsection{JIRA Project Board}
	% Required screenshot(s) for your JIRA project board.
	\noindent \textit{[Insert screenshots of 'Summary', 'timeline', and 'backlog' pages here]}
	
	\subsection{Individual Reflection Essays}
	% Each group member's individual reflection essay.
	\subsubsection{Luxiao Cao}
	[Insert 500-word reflection here]
	\subsubsection{Akshansh Rajora}
	[Insert 500-word reflection here]
	\subsubsection{Fuyu Cao}
	[Insert 500-word reflection here]
	\subsubsection{Tianwei Yu}
	[Insert 500-word reflection here]
	\subsubsection{Detai Dong}
	[Insert 500-word reflection here]
	
	\subsection{Team Sprint Reports}
	% All team sprint reports in chronological order.
	[Insert chronological sprint reports here]
	
	\subsection{Complete Source Code}
	% The complete DApp Solidity smart contract code and testing codes. All in text format.
	\subsubsection{Smart Contracts}
	\begin{verbatim}
		// Insert complete Solidity code here
	\end{verbatim}
	
	\subsubsection{Frontend Core Files}
	\begin{verbatim}
		// Insert modified/relevant frontend files (e.g., page.js) here
	\end{verbatim}
	
\end{document}
